\section{Feedback}

\begin{frame}{Feedback}
    \small
    \textcolor{red_unipd}{\Large Midterm Questions (1/2)}
    \vspace{0.3em}

    \scriptsize
    \begin{enumerate}
        \item \textbf{What papers are you referencing?}
        \item What is Akash's physical theory/model?
        \item \textbf{How do we generate the dataset?}
        \item \textbf{Why are you choosing Burger's equation?}
        \item What are we trying to simulate?
        \item \textbf{What are the current limitations in related papers? What approaches are they trying?}
        \item Related work - Any similar simulation?
        \item What about accuracy benchmarks? (Note: Increase speed, decrease accuracy trade-off)
        \item What kind of initial conditions are you using?
        \item What distributions are you using?
    \end{enumerate}
\end{frame}

\begin{frame}{Feedback}
    \small
    \textcolor{red_unipd}{\Large Midterm Questions (2/2)}
    \vspace{0.3em}

    \scriptsize
    \begin{enumerate}
        \setcounter{enumi}{10}
        \item Use the same notation for everything - Bigger figures
        \item Compare your architecture vs NFTM
        \item What is the main conclusion per model?
        \item Show loss graphs
        \item What is the absolute training error? Is it happening during training?
        \item Akash - Show training error, should be almost perfect
        \item Show the physical analogy
        \item For the context window - 5 timesteps \& at least 2-3 viscosity values for training
        \item Related work always - What is the current state of the art?
    \end{enumerate}
\end{frame}

\begin{frame}{Feedback}
    \small
    \textcolor{red_unipd}{\Large Answers to Key Questions (1/2)}
    \vspace{0.3em}

    \scriptsize
    \textbf{1. What papers are you referencing?}
    \begin{itemize}
        \item Neural Turing Machines (Graves, Wayne, Danihelka, 2014)
        \item Physics-Informed Neural Networks (Raissi, Perdikaris, Karniadakis, 2019)
        \item Fourier Neural Operator (Li et al., 2021)
        \item Neural Operator (Kovachki et al., 2023)
        \item Spectral Bias work (Rahaman et al., 2019)
        \item Neural Field Turing Machine (Malhotra \& Seghouani, 2025)
    \end{itemize}

    \vspace{0.5em}
    \textbf{3. How do we generate the dataset?}
    \begin{itemize}
        \item Uniform mesh: N=128 spatial points, 256 time steps, $\Delta t=5\times10^{-3}$
        \item Sample initial conditions (shock, rarefaction, sinusoidal, smoothed random)
        \item Conservative flux integration (Rusanov) with diffusion term
        \item CFL-safe time steps with stability checks; reject unstable trajectories
        \item Store as .npz files with fields U, x, t, $\nu$, $\Delta x$, $\Delta t$
    \end{itemize}
\end{frame}

\begin{frame}{Feedback}
    \small
    \textcolor{red_unipd}{\Large Answers to Key Questions (2/2)}
    \vspace{0.3em}

    \scriptsize
    \textbf{4. Why are you choosing Burger's equation?}
    \begin{itemize}
        \item Simple enough to learn with limited data/resources, yet rich enough to exhibit complex dynamics (shock formation, nonlinear advection)
        \item 1D simplified version of Navier-Stokes that captures key fluid dynamics features
    \end{itemize}

    \vspace{0.5em}
    \textbf{6. What are the current limitations in related papers?}
    \begin{itemize}
        \item \textbf{PINNs}: Spectral bias toward low frequencies, struggle with shocks, slow convergence ($>10^5$ iterations)
        \item \textbf{NTMs}: Discrete memory unsuitable for continuous spatial fields, no spatial inductive biases
        \item \textbf{FNO/Neural Operators}: Serve as reference benchmarks
    \end{itemize}
\end{frame}
